\subsubsection{Global Reach}
\label{sec:glo-reach}
In this subsection, we estimate the lower bound of the global reach of $\varphi(\MM)$. The global reach constrains the overall shape of the manifold by controlling the Euclidean distance to ''separate'' different parts of the manifold.


The proof is organized by considering three different cases depending on different ranges of the geodesic distance $d(p,q)$. When $d(p,q)$ is small, a local linear approximation excludes the global case directly (\cref{lem:global_case}); when $d(p,q)$ is large, heat-kernel estimates give a positive lower bound on $|\varphi(p)-\varphi(q)|$ (\cref{global condition}); the remaining is closed either by the condition $(*)$ or by a refined heat-kernel analysis if $(*)$ fails.

Under Assumptions \ref{as:reach}, when $p$ is close to $q$, the geodesic
distance $d(p,q)$ can be controlled by the Euclidean distance $|p-q|$. 

\begin{lemma}\cite[Prop. 1]{li2019geodesic}
	\label[lemma]{lem:geodis_eucdis}
	If $d(p,q)=s$, then
	\begin{align}
		s-\frac{s^3}{24 r_0^2}\leq |p-q| \leq s,
	\end{align}
	where $\frac{1}{r_0}=\sup
	\left\{|\gamma^{\prime\prime}(s)|\right\}$ and $\gamma$ varies among
	all geodesics on $M$ in arc length parameter.
\end{lemma}

We consider the left inequality of \cref{lem:geodis_eucdis}, it is easy to see if $s\leq 2\sqrt{2}r_0$, then
\begin{align}
	\frac{2}{3}s\leq s-\frac{s^3}{24r_0^2} \leq |p-q|,
\end{align}
which implies 
\begin{align}
	|p-q|\leq d(p,q) \leq \frac{3}{2}|p-q|.
\end{align}

We can verify easily that $1/r_0=\sup
\left\{|\gamma^{\prime\prime}(s)|\right\}=1/\tau_{l}$, and this local linear approximation can exclude the global reach case when $d(p,q)< 2\sqrt{2}\tau_{l}$, and $\frac{3}{2}$ plays an important role here, it will lead to the contradiction to the global reach case.

\begin{lemma}
	\label[lemma]{lem:global_case}
	If $d(p,q)\leq s_0$ where $s_0=2\sqrt{2}\tau_{l}$, then $p,q$ cannot satisfy global reach case.
\end{lemma}

\begin{proof}[Proof of \cref{lem:global_case}]
	\label{proof:global_case}
	If not, we assume there exists $p,q \in M$ such that $d(p,q)\leq s_0$ and $p,q$ satisfy global reach case, i.e. 
	\begin{align}
		
&\frac{p+q}{2} \in Med(M), \\
 \pi(\frac{p+q}{2})=p,q &\quad \text{and} \quad d_E(p,M)=d_E(q,M)=|p-q|/2 \label{proj_med}
	\end{align}
Since $d(p,q)\leq s_0$, we have $d(p,q)\leq \frac{3}{2}|p,q|$. 
	
	
	\begin{figure}[H]
		\centering
		\begin{tikzpicture}[scale=2]
			
			% 画圆
			\draw[dashed] (0,0) circle(1cm);
			
			% 画直径
			\draw[dashed] (-0.6,-0.8) -- (0.6,0.8);
			
			% 画穿过直径与圆交点的曲线
			\draw plot[smooth, tension=0.7] coordinates {(-3,-1.5) (-1,-1.2) (-0.6,-0.8) (-0.6,0.2) (0.6,0.8) (1,1) (3,1.5)};
			\node at (-0.6,-0.8) [left] {$p$};
			\node at (0.6,0.8) [above] {$q$};
			\node at (0,0) [below right] {$\frac{p+q}{2}$};
			\node at (-0.15,0.5) [above left] {$\gamma(a)=\pi(\frac{p+q}{2})$};
			\filldraw[red] (-0.6,-0.8) circle (1pt);
			\filldraw[red] (0.6,0.8) circle (1pt);
			\filldraw[red] (0,0) circle (1pt);
			\filldraw[red] (-0.21,0.47) circle (1pt);
			\draw[dashed] (0, 0) -- (-0.21,0.47);
		\end{tikzpicture}
		\caption{Local Estimate}
	\end{figure}
	
	The inequality above implies the geodesic $\gamma(s)$ connecting $p$ and
	$q$ must be partly inside the ball
	$B(\frac{p+q}{2},\frac{|p-q|
	}{2})$, since $\overline{pq}$ is the diameter of this ball and
	\begin{align}
		|\gamma(s)|=d(p,q) \leq \frac{3}{2} |p-q| < \frac{\pi}{2}|p-q|,
	\end{align}
	which implies there exist $a \in (0,s)$ such that
	\begin{align}
		|\gamma(a)-\frac{p+q}{2}| < \frac{|p-q|}{2},
	\end{align}
	which contradicts with $\varphi(p),\varphi(q)$ satisfy global reach case \cref{proj_med}.
\end{proof}


Since $\varphi$ is almost isometry, i.e. $| \|d\varphi\|-1| < \epsilon$, thus $d(p,q)\leq s/(1+\epsilon)$ implies $d(\varphi(p),\varphi(q)) \leq s$. Since we have proved that the local reach of $\varphi(M)$ admits a uniform lower bound $\tau_{l,\varphi(\MM)}$ depending only on the geometric parameters of $\MM$, then we choose $s_0 = 2\sqrt{2}\tau_{l,\varphi(\MM)}$, then $\varphi(p), \varphi(q)$ cannot satisfy global reach case for any $\varphi(M) \in \varphi(\MM)$ Consequently, only when $d(p,q)>\frac{s_0}{1+\epsilon}$, $\varphi(p),\varphi(q)$ can satisfy global reach case. 


And in the following, we will claim that if the geodesic distance of two points $p,q \in M$ is large enough, then the Euclidean distance of $\varphi(p), \varphi(q) \in \varphi(M)$ is bounded from below, i.e. 
\begin{lemma}
	\label[lemma]{global condition}
	For large enough $s_1>0$, there exists $r_1>0$ such that
	\begin{align}
		d(p,q)> s_1 \implies |\varphi(p)-\varphi(q)|>r_1.
	\end{align}
	We call this the {\em global reach condition}.
\end{lemma}




To bound the Euclidean distance from below, we need to bound the heat kernel $K(t,p,q)$ from above and $K(t,p,p)$ from below to control this Euclidean distance. For upper bound, we have
\begin{lemma}
	\label{thm:upperbound}
	Let M be a complete Riemannian manifold of dimension $d$ with Ricci curvature is greater than $-\kappa(d-1)$ for some $\kappa\geq 0$, then heat kernel satisfies:
	\begin{align}
		K(t,p,q)\leq \frac{C_1(d)}{t^{d/2}}\exp\left(C_2(d)\kappa t - \frac{2d(p,q)^2}{9t}\right)
	\end{align}.
\end{lemma}

This lemma follows from \cref{thm:upper_app} of Li-Yau and \cref{injvol} of Croke with $\alpha_1=\frac{3}{2}$, $\alpha_2=\frac{1}{2}$ and $t_0\leq \iota^2/4$, and \cref{thm:upper_app}, \cref{injvol} are in the appendix.


For lower bound, we only need on-diagonal lower estimate:
\begin{lemma}\cite[Thm. 3.1]{wang1997sharp}
	\label{thm:lowerbound}
	Let M be a complete Riemannian manifold of dimension $d$ with Ricci curvature is greater than $-\kappa(d-1)$ for some $\kappa\geq 0$. For any $t>0$ and $p\in M$, we have
	\begin{align}
		K(t,p,p) \geq (4\pi t)^{-d/2} \exp\left(-\frac{\beta^2}{4}t-\frac{2\sqrt{3d}\beta}{3}\sqrt{t}\right).
	\end{align}
	where $\beta=\sqrt{\kappa}(d-1)$.
\end{lemma}
This follows \cref{thm:lower_app} of Wang with $\sigma^2=\frac{3\beta^2}{8d}$ and $p=q$.



\begin{proof}[Proof of \cref{global condition}]
	To show this, we will compute the distance between $\varphi(p)$ and $\varphi(q)$ in $\mathbb{R}^m$ directly and then estimate it using geodesic distance of $p,q$.
	
	For any $p \in M$, $\varphi(p)=(t_0)^\frac{d+2}{4}\sqrt{2}(4\pi)^\frac{d}{4}\left(e^{-\lambda_1 t_0/2}e_1(p),\cdots,e^{-\lambda_m t_0/2}e_m(p)\right) \in \mathbb{R}^m$, thus the Euclidean distance between $\varphi(p)$ and $\varphi(q)$ is
	\begin{align}
		|\varphi(p)-\varphi(q)|&=(t_0)^\frac{d+2}{4}\sqrt{2}(4\pi)^\frac{d}{4} \sqrt{\sum_{i=1}^me^{-\lambda_i t_0}(e_i(p)-e_i(q))^2}\\
		&=(t_0)^\frac{d+2}{4}\sqrt{2}(4\pi)^\frac{d}{4}\sqrt{\sum_{i=1}^me^{-\lambda_i t_0}(e_i(p)^2+e_i(q)^2-2e_i(p)e_i(q))} \\
		&=(t_0)^\frac{d+2}{4}\sqrt{2}(4\pi)^\frac{d}{4}\sqrt{\sum_{i=0}^m(e^{-\lambda_i t_0}e_i(p)^2+e^{-\lambda_i t_0}e_i(q)^2-2e^{-\lambda_i t_0}e_i(p)e_i(q))},
	\end{align}
	the last equality holds since $e_0$ is a constant function.
	
	
	We have the representation of heat kernel
	\begin{align}
		K(t,p,q)=\sum_{i=0}^{\infty} e^{-\lambda_i t} e_i(p) e_i(q),
	\end{align}
	and the truncated heat kernel $K_m$ is
	\begin{align}
		K_m(t,p,q)=\sum_{i=0}^{m} e^{-\lambda_i t} e_i(p) e_i(q),
	\end{align}
	and we have $\|K_m(t_0,p,\cdot)-K(t_0,p,\cdot)\|_{\infty} \leq \epsilon^\prime$,
	thus 
	\begin{align}
		|\varphi(p)-\varphi(q)|&\geq (t_0)^\frac{d+2}{4}\sqrt{2}(4\pi)^\frac{d}{4} \sqrt{K_m(t_0,p,p)+K_m(t_0,q,q)-2K_m(t_0,p,q)}\\
		&\geq (t_0)^\frac{d+2}{4}\sqrt{2}(4\pi)^\frac{d}{4} \sqrt{K(t_0,p,p)+K(t_0,q,q)-2 K(t_0,p,q)-4\epsilon^\prime}.
	\end{align}
	
	Combining two inequities in \cref{thm:upperbound} and \cref{thm:lowerbound}, we have
	\begin{align}
		|\varphi(p)&-\varphi(q)|\geq (t_0)^\frac{d+2}{4}2(4\pi)^\frac{d}{4}  \nonumber \\ \times &\sqrt{(4\pi t_0)^{-\frac{d}{2}} \exp\left(-\frac{\beta^2t_0}{4}-\frac{2\sqrt{3dt_0}\beta}{3}\right) - \frac{C_1(d)}{t_0^{d/2}}\exp\left(C_2(d)\kappa t_0 - \frac{2d(p,q)^2}{9t_0}\right) - 2\epsilon^\prime}. \label{sqrt_dis}
	\end{align}
	In our case, $t_0, m$ are fixed when embedding $M$ into $\mathbb{R}^m$. The first two items in the square root
	are only depends on the geometric properties of $\MM$, and $\epsilon^\prime$ is chosen when we perform DM, we can choose appropriate $\epsilon^{\prime}$ based on our geometric setting such that 
	\begin{align}
		2\epsilon^{\prime}<(4\pi t_0)^{-\frac{d}{2}} \exp\left(-\frac{\beta^2t_0}{4}-\frac{2\sqrt{3dt_0}\beta}{3}\right)
	\end{align}
	
	Therefore, for large enough $s_1$ with $d(p,q)\geq s_1$, we have $|\varphi(p)-\varphi(q)|\geq r_1$, where $r_1, s_1$ depend on our
	setting of $\MM$.
\end{proof}

The proof above only claimed the existence of $s_1$ and $r_1$. For the uniformity claim we now make them explicit, choosing $\epsilon^\prime$, $s_1$ and $r_1$ so that they depend only on the geometric parameters of $\MM$, we consider the item under square root in \cref{sqrt_dis} which is positive, that is
\begin{align}
	&(4\pi t_0)^{-\frac{d}{2}} \exp\left(-\frac{\beta^2t_0}{4}-\frac{2\sqrt{3dt_0}\beta}{3}\right) - \frac{C_1(d)}{t_0^{d/2}}\exp\left(C_2(d)\kappa t_0 - \frac{2s_1^2}{9t_0}\right) - 2\epsilon^\prime > 0 \\
	\iff&  \exp\left(C_2(d)\kappa t_0 - \frac{2s_1^2}{9t_0}\right)< \frac{(4\pi)^{-\frac{d}{2}}}{C_1(d)} \exp\left(-\frac{\beta^2t_0}{4}-\frac{2\sqrt{3dt_0}\beta}{3}\right) - \frac{2(t_0)^\frac{d}{2}\epsilon^\prime}{C_1(d)}.
\end{align}

We denote the right side $\frac{(4\pi)^{-\frac{d}{2}}}{C_1(d)} \exp\left(-\frac{\beta^2t_0}{4}-\frac{2\sqrt{3dt_0}\beta}{3}\right) - \frac{2(t_0)^\frac{d}{2}\epsilon^\prime}{C_1(d)}$ as $F$, we need it be a fixed constant, thus we can select 
\begin{align}
	\label{sele eprime}
	\epsilon^\prime=\frac{(4\pi t_0)^{-\frac{d}{2}}}{8}\exp\left(-\frac{\beta^2t_0}{4}-\frac{2\sqrt{3dt_0}\beta}{3}\right),
\end{align}
which implies
\begin{align}
	F= \frac{3(4\pi)^{-\frac{d}{2}}}{4C_1(d)} \exp\left(-\frac{\beta^2t_0}{4}-\frac{2\sqrt{3dt_0}\beta}{3}\right),
\end{align}
Therefore, $\epsilon^\prime$ is only depending on the geometric properties of $\MM$. Choosing $\epsilon^\prime$ will affect $m=\max\left\{N_0(d,\kappa,\iota,\epsilon,V,\frac{t_0}{2}),N_1(d,\kappa,\iota,V,\epsilon^\prime,t_0)\right\}+1$ and is independent with choice of $t_0=t_0(d,\kappa,\iota,\epsilon)$ when perform DM. We need set appropriate $\epsilon^\prime$ based on geometric setting and derive $m$ and $t_0$ at the beginning, but until now we know how to choose them.

Then we ask $\exp\left(C_2(d)\kappa t_0 - \frac{2s_1^2}{9t_0}\right) < \frac{2}{3}F$, which is equivalent to
\begin{align}
	s_1^2>\frac{9t_0}{2}\left(C_2(d)\kappa t_0 + \frac{\beta^2 t_0}{4}+\frac{2\sqrt{3d t_0}\beta}{3}+ \log\left(2(4\pi)^\frac{d}{2}C_1(d)\right)\right), \label{smallests1}
\end{align}
Intuitively, $t_0$ is small, so $s_1$ is not very large.

When $s_1$ is the square root of \cref{smallests1} and  $\epsilon^\prime=\frac{(4\pi t_0)^{-\frac{d}{2}}}{8}\exp\left(-\frac{\beta^2t_0}{4}-\frac{2\sqrt{3dt_0}\beta}{3}\right)$, we have 
\begin{align}
	r=r_1=\sqrt{t_0} \exp\left(-\frac{\beta^2t_0}{8}-\frac{\sqrt{3dt_0}\beta}{3}\right)
\end{align}

We now split $d(p,q)$ into three cases according to the range, which together cover all pairs $p,q\in M$:
\begin{enumerate}
	\item [(i)] $d(p,q)\leq s_0/(1+\epsilon)$. Since $|\|d\varphi\|-1|<\epsilon$, we have $d(\varphi(p),\varphi(q))\leq s_0$, then by \cref{lem:global_case}, $\varphi(p), \varphi(q)$ cannot be the global reach case.
	\item [(ii)] $d(p,q)>s_1$. It implies $|\varphi(p)-\varphi(q)|> r_1$ by \cref{global condition}, we have the reach is bounded from the below.
	\item [(iii)] $s_0/(1+\epsilon) < d(p,q) < s_1$. If there is no gap between $s_0/1+\epsilon$ and $s_1$, that is
	\begin{align}
		s_0^2>\frac{9(1+\epsilon)^2t_0}{2}\left(C_2(d)\kappa t_0 + \frac{\beta^2 t_0}{4}+\frac{2\sqrt{3d t_0}\beta}{3}+\log\left(2(4\pi)^\frac{d}{2}C_1(d)\right)\right).
	\end{align}
	Combining $s_0=2\sqrt{2}r_0$ and $r_0 =\tau_{l,\varphi(\MM)}$. The condition making $s_0/1+\epsilon \geq s_1$ is 
	\begin{align}
		\label{star}
		8\tau^2_{l,\varphi(\MM)} &\geq 8/\left(d^2(t_0)^\frac{d+2}{2}2(4\pi)^\frac{d}{2} C_2^2\sum_{k=1}^{m}e^{-\lambda_k t_0} (1+\lambda_k)^{2+\frac{d}{2}} \right) \nonumber \\
		&\geq \frac{9(1+\epsilon)^2t_0}{2}\left(C_2(d)\kappa t_0 + \frac{\beta^2 t_0}{4}+\frac{2\sqrt{3d t_0}\beta}{3}+ \log\left(2(4\pi)^\frac{d}{2}C_1(d)\right)\right), \tag{*}
	\end{align}
	which is an assumption only based on geometric properties since $t_0=t_0(d,\kappa,\iota,\epsilon)$, $m=\max\left\{N_0,N_1\right\}+1$ and so on. If $(*)$ fails, there is an uncovered interval $(s_0/(1+\epsilon),\,s_1)$, we need to verify $d(p,q)$ in it. To do this, we need a tighter lower bound on $s_0$, and a short time expansion of the heat kernel that controls $K(t,p,p)-K(t,p,q)$ in this interval.
\end{enumerate}


Now we consider (iii) when there is a gap. First we estimate $\tau_{l,\varphi(M)}$ in more details. From now, we ignore some constants and merge them as $C$ to make our computation clearly. We have $\| II\|_{op} \leq C(t_0)^{\frac{d+2}{4}} \sqrt{\sum_{k=1}^{m}e^{-\lambda_k t_0} (1+\lambda_k)^{2+\frac{d}{2}}}$, combining the estimate for eigenvalues
\begin{align}
C_1^{1+\text{diam}\sqrt{\kappa}} \text{diam}^{-2} k^{2/d} \leq \lambda_k \leq C+C^\prime k^{2/d},
\end{align}
we can derive
\begin{align}
\|II\|_{op} \leq C (t_0)^{\frac{d+2}{4}}\sqrt{\sum_{k=1}^{m}e^{-k^{2/d} t_0} k^{1+\frac{4}{d}}} < C (t_0)^{\frac{d+2}{4}}\sqrt{\sum_{k=1}^{\infty}e^{-k^{2/d} t_0} k^{1+\frac{4}{d}}}.
\end{align} 
Now we estimate $\sum_{k=1}^{\infty}e^{-k^{2/d} t_0} k^{1+\frac{4}{d}}$ by integral. We have 
\begin{align}
	\sum_{k=1}^{\infty}e^{-k^{2/d} t_0} k^{1+\frac{4}{d}} \leq  \int_{0}^{\infty} e^{-x^{2/d} t_0} x^{1+\frac{4}{d}}dx + C
\end{align}
let $u = x^{2/d} t_0$, then $x^{1+\frac{4}{d}} = (\frac{u}{t_0})^{2+\frac{d}{2}}$ and $dx = \frac{u^{\frac{d}{2}-1}d}{2 t_0^{d/2}} du$, then 
\begin{align}
	\int_{0}^{\infty} e^{-x^{2/d} t_0} x^{1+\frac{4}{d}}dx = \frac{d}{2t_0^{2+d}}\int_{0}^{\infty} e^{-u} u^{d+1} du = \frac{d}{2t_0^{2+d}}\Gamma(d+2)
\end{align}

Therefore, $\|II\|_{op}^2 \leq Ct_0^{\frac{d+2}{2}}/t_0^{d+2} = Ct_0^{-(1+\frac{d}{2})}$, and $s_0^2 = 8 \tau_{l,\varphi(M)} \geq C t_0^{1+\frac{d}{2}}$.

Now we try to estimate the diffusion distance carefully using the asymptotic expression, which is due to Varadhan:
\begin{align}
	\lim_{t \to 0} -4t \log K(t,p,q) = d(p,q)^2
\end{align}
and thus
\begin{align}
	\lim_{t \to 0} -4t \log \frac{K(t,p,q)}{K(t,p,p)} = d(p,q)^2
\end{align}
which implies $K(t,p,q) \approx K(t,p,p) e^{-\frac{d(p,q)^2}{4t}}$ for small enough $t$. Then 
\begin{align}
	K(t,p,p) - K(t,p,q) \approx (1 - e^{-\frac{d(p,q)^2}{4t}}) K(t,p,p).
\end{align}

To make it more acurate, we have 
\begin{align}
	\label{eq:heatkernelcoef}
	\left|K(t,p,q)- \frac{e^{-d(p,q)^2/4t}}{(4\pi t)^{d/2}}\sum_{j=0}^{v} t^j \frac{\Phi_j(p,q)}{j !}\right| \leq C(v) t^{v+1} \frac{e^{-d(p,q)^2/4t}}{(4\pi t)^{d/2}}
\end{align}
holds for any $v \in \N$, all $(x,y)$ not in the cut locus and $t$ is less than one specified $T$, and $\Phi_j(p,q)$ are the heat kernel coefficients \cite{ludewig2019strong}.

Then we have
\begin{align}
	\frac{K(t,p,q)}{K(t,p,p)} \leq \frac{\frac{e^{-d(p,q)^2/4t}}{(4\pi t)^{d/2}} (\sum_{j=0}^{v} t^j \frac{\Phi_j(p,q)}{j !} + C(v) t^{v+1})}{\frac{1}{(4\pi t)^{d/2}} (\sum_{j=0}^{v} t^j \frac{\Phi_j(p,p)}{j !} - C(v) t^{v+1})}
\end{align}
thus
\begin{align}
	K(t,p,p) - K(t,p,q) \geq \left(1- \frac{ \sum_{j=0}^{v} t^j \frac{\Phi_j(p,q)}{j !} + C(v) t^{v+1}}{\sum_{j=0}^{v} t^j \frac{\Phi_j(p,p)}{j !} - C(v) t^{v+1}}e^{-d(p,q)^2/4t}\right) K(t,p,p).
\end{align}

We assume that we can find $t$ and $v$ such that $\frac{ \sum_{j=0}^{v} t^j \frac{\Phi_j(p,q)}{j !} + C(v) t^{v+1}}{\sum_{j=0}^{v} t^j \frac{\Phi_j(p,p)}{j !} - C(v) t^{v+1}} \leq e^{d(p,q)^2/8t}$ when $d(p,q)^2 \in (s_0^2,s_1^2)$. We also have
\begin{align}
	K(t,p,p) \geq (4\pi t)^{-d/2} \exp\left(-\frac{\beta^2}{4}t-\frac{2\sqrt{3d}\beta}{3}\sqrt{t}\right).
\end{align}
thus 
\begin{align}
	K(t,p,p) - K(t,p,q) \geq \frac{1 - e^{-\frac{d(p,q)^2}{8t}}}{(4\pi t)^{d/2}} \exp\left(-\frac{\beta^2}{4}t-\frac{2\sqrt{3d}\beta}{3}\sqrt{t}\right).
\end{align}
when $t$ is small, $\exp\left(-\frac{\beta^2}{4}t-\frac{2\sqrt{3d}\beta}{3}\sqrt{t}\right)$ is obviously bounded from below, we denote the lower bound as $C_l$. The remaining part increases monotonically with $d(p,q)^2$ in $(s_0^2, s_1^2)$, so 
\begin{align}
	K(t,p,p)- K(t,p,q) \geq \frac{1 - e^{-\frac{C}{8} t^{d/2}}}{(4\pi t)^{d/2}} C_l
\end{align}
which is converge to a positive constant as $t \to 0$. This implies $K(t,p,p) - K(t,p,q)$ has a positive lower bound with an appropriate $t$, also $|\varphi(p) - \varphi(q)|$ has the lower bound. So when $d(p,q)^2 \geq s_0^2 = Ct^{1+\frac{d}{2}} $, we can always find the lower bound of $|\varphi(p)- \varphi(q)|$. Hence the global reach are bounded from below if we choose appropriate $t$.



\iffalse

When $\epsilon\rightarrow 0$, $\epsilon^\prime \rightarrow 0 $, the inequality is
\begin{align}
	s_0^2>\frac{9t_0}{2}\left(C_2(d)\kappa t_0 + \frac{\beta^2 t_0}{4}+\frac{2\sqrt{3d t_0}\beta}{3}+ \frac{d}{2}\log(4\pi)+\log C_1(d)\right).
\end{align}

This inequality is the necessary condition, which describes the requirement in the case with no error. If this inequality holds, we can select small enough $\epsilon$ and $\epsilon^\prime$ when perform DM, the global reach condition holds.

\fi


Therefore, the local reach and global reach both have the uniform lower bound, thus $\tau_{\varphi(M)}\geq \tau_{\varphi(\MM)}>0$ for some fixed constant.
